\subsection{Fattoriali, coefficienti binomiali e Bernoulli}
\label{sec:esercizi-fondamenti-fattoriali-binomiali-bernoulli}

\begin{exercise}[Riduzione di fattoriali]
\label{ex:fondamenti-semplificazione-fattoriali}
\difficoltaesercizio{0}
Semplificare, per ogni $n\in\N$, l'espressione
\[
  \frac{(n+3)!}{n!\,(n+1)(n+3)}.
\]
\rimandoteoria{def:fattoriale}
\end{exercise}
\begin{proof}[Soluzione dell'esercizio]
Per la definizione di fattoriale,
\[
\begin{aligned}
  \frac{(n+3)!}{n!\,(n+1)(n+3)}
    &=\frac{(n+3)(n+2)(n+1)n!}{n!\,(n+1)(n+3)}\\
    &=n+2.
\end{aligned}
\]
\end{proof}

\begin{exercise}[Coefficiente binomiale elementare]
\label{ex:fondamenti-coefficiente-binomiale-due}
\difficoltaesercizio{1}
Dimostrare che, per ogni $n\ge2$,
\[
  \binom{n}{2}=\frac{n(n-1)}{2}.
\]
\rimandoteoria{def:coefficiente-binomiale}
\end{exercise}
\begin{proof}[Soluzione dell'esercizio]
Per la definizione di coefficiente binomiale,
\[
\begin{aligned}
  \binom{n}{2}
    &=\frac{n!}{2!\,(n-2)!}\\
    &=\frac{n(n-1)(n-2)!}{2(n-2)!}\\
    &=\frac{n(n-1)}{2}.
\end{aligned}
\]
\end{proof}

\begin{exercise}[Somma di una riga del triangolo di Pascal]
\label{ex:fondamenti-somma-coefficienti-binomiali}
\difficoltaesercizio{1}
Calcolare $\sum_{k=0}^{4}\binom{4}{k}$ e successivamente dimostrare
che, per ogni $n\in\N$,
\[
  \sum_{k=0}^{n}\binom{n}{k}=2^n.
\]
\rimandoteoria{thm:binomio-newton}
\end{exercise}
\begin{proof}[Soluzione dell'esercizio]
Per $n=4$,
\[
  \sum_{k=0}^{4}\binom{4}{k}=1+4+6+4+1=16.
\]
Per il caso generale si applica il binomio di Newton con $a=b=1$:
\[
  \sum_{k=0}^{n}\binom{n}{k}
    =\sum_{k=0}^{n}\binom{n}{k}1^k1^{n-k}
    =(1+1)^n=2^n.
\]
\end{proof}

\begin{exercise}[Identità tra coefficienti binomiali]
\label{ex:fondamenti-identita-coefficienti-binomiali}
\difficoltaesercizio{1}
Dimostrare che, per $n\ge k\ge1$,
\[
  k\binom{n}{k}=n\binom{n-1}{k-1}.
\]
\rimandoteoria{def:coefficiente-binomiale}
\end{exercise}
\begin{proof}[Soluzione dell'esercizio]
Dalla definizione segue
\[
\begin{aligned}
  k\binom{n}{k}
    &=k\frac{n!}{k!(n-k)!}\\
    &=\frac{n!}{(k-1)!(n-k)!}\\
    &=n\frac{(n-1)!}{(k-1)!(n-k)!}\\
    &=n\binom{n-1}{k-1}.
\end{aligned}
\]
\end{proof}

\begin{exercise}[Stima mediante la disuguaglianza di Bernoulli]
\label{ex:fondamenti-bernoulli-stima}
\difficoltaesercizio{1}
Dimostrare che, per ogni $n\ge1$,
\[
  \left(1+\frac1n\right)^n\ge2.
\]
Determinare quando vale l'uguaglianza.
\rimandoteoria{thm:disuguaglianza-bernoulli}
\end{exercise}
\begin{proof}[Soluzione dell'esercizio]
Applicando la disuguaglianza di Bernoulli con $x=1/n>0$, si ha
\[
  \left(1+\frac1n\right)^n
  \ge 1+n\cdot\frac1n=2,
\]
come richiesto. L'uguaglianza si verifica per $n=1$; per $n\ge2$
lo sviluppo binomiale contiene, oltre ai termini $1$ e $n/n$,
anche termini positivi, quindi la disuguaglianza è stretta.
\end{proof}

\clearpage

\begin{exercise}[Confronto tra fattoriali e potenze]
\label{ex:fondamenti-fattoriali-potenze}
\difficoltaesercizio{1}
Dimostrare per induzione che, per ogni \(n\ge5\),
\[
  n!\ge3^{n-1}.
\]
Verificare inoltre perché l'enunciato sarebbe falso per \(n\ge4\).
\rimandoteoria{thm:principio-induzione}
\end{exercise}
\begin{proof}[Soluzione dell'esercizio]
Per \(n=0,1,2,3,4\) si ha
\[
\begin{aligned}
  0!&=1>\frac13=3^{-1},\\
  1!&=1=3^0,\\
  2!&=2<3=3^1,\\
  3!&=6<9=3^2,\\
  4!&=24<27=3^3.il
\end{aligned}
\]
Pertanto la disuguaglianza vale per \(n=0,1\),
ma non per \(n=2,3,4\). In particolare,
l'enunciato con \(n\ge4\) è falso.

Si pone
\[
  A:=\left\{n\in\N:n\ge5\ \text{e}\
  n!\ge3^{n-1}\right\}.
\]

\emph{Passo base.} Per \(n=5\),
\[
  5!=120\ge81=3^4,
\]
quindi \(5\in A\).

\emph{Passo induttivo.} Sia \(n\ge5\) e si supponga \(n\in A\). Allora
\[
\begin{aligned}
  (n+1)!
    &=(n+1)n!\\
    &\ge(n+1)3^{n-1}\\
    &\ge3\cdot3^{n-1}\\
    &=3^n,
\end{aligned}
\]
poiché \(n+1\ge6>3\).

Pertanto \(n\in A\Rightarrow n+1\in A\) e la tesi segue per induzione.
\end{proof}
